% Notes of the Montreal seminar, 1970 (cote 129, typed part): archive cote 129, site page 39 (PDF page 40).
% typescript. Draft compared: none.
% AI-transcribed by Claude Opus 5.5 (claude-opus-5-5, adaptive thinking, effort high), Message Batches API, batch mtl129_full (msgbatch_01AGnDKHwy6ATCM7ehUa6263), 2026-11-13; not yet checked by eye.
\scanpage{129/39}
\mspage{\hand{II.11}}

Si $k(s)$ est de caractéristique $p$ et si $G$ est annulé
par une puissance de $p$ , $G^{\mathrm{mult}}$ est infinitésimal (2.2) et on a donc
des inclusions
\[ 0 \subset G^{\mathrm{mult}} \subset G^{0} \subset G \]
\margin{X} où $G/G^{0} = G^{\mathrm{et}}$ est étale, $G^{0}/G^{\mathrm{mult}}$ est bi-infinitésimal\hand{,}\deleted{(}car c'est
un quotient d'un groupe infinitésimal et son dual s'identifie à $((G^{0})^{*})^{0}$ .

Cette suite de composition de $G$ est fonctorielle en $G$ ,
et on vérifie qu'elle commute aux produits finis.

3.2. Si $S$ est le \uline{spectre d'un corps parfait de caractéristique} $p$ ,
la projection canonique $G \to G^{\mathrm{et}}$ induit un isomorphisme de $G_{\mathrm{red}}$ sur
$G^{\mathrm{et}}$ et la suite exacte $(*)$ splitte canoniquement (SGA 3, XVII, 1.6),
et on a donc
\[ G = G^{0} \times G^{\mathrm{et}} . \]

De même la suite exacte
\[ 0 \to G^{\mathrm{mult}} \to G^{0} \to G^{0}/G^{\mathrm{mult}} \to 0 \]
est splittée car la suite duale l'est. On obtient donc une décomposition
canonique, fonctorielle en $G$ :
\[ G = G^{\mathrm{mult}} \times G^{\mathrm{bi}} \times G^{\mathrm{et}} \]
où $G^{\mathrm{mult}}$, $G^{\mathrm{bi}} = G^{0}/G^{\mathrm{mult}}$ , $G^{\mathrm{et}}$ sont respectivement de type multiplica-%
tif, bi-infinitésimal et étale. La partie unipotente de $G$ est alors
$G^{\mathrm{bi}} \times G^{\mathrm{et}}$ et la partie infinitésimale est $G^{\mathrm{mult}} \times G^{\mathrm{bi}}$ .

\mspage{\hand{39}}
